\documentclass[UTF8, 11pt]{ctexart}

\usepackage[margin=1in]{geometry}
\usepackage{amsmath, amssymb}
\usepackage{booktabs}
\usepackage{graphicx}
\usepackage{hyperref}
\usepackage{url}
\usepackage{enumitem}
\setlist{leftmargin=*}
\setlength{\parskip}{0.35em}
\setlength{\parindent}{0pt}

\title{\textbf{阶段 4：实验评估}\\
运行时间、可扩展性与基线对比}
\author{小组成员：姓名 1，姓名 2，姓名 3}
\date{CS240：算法设计与分析，2026 年春季学期}

\begin{document}

\maketitle

\section{评估目标}

本阶段的目标是评估实现的算法是否支持项目核心观点：经典图划分算法能够在保持
工作负载均衡的同时，降低分布式优化中的通信成本。实验评估对应课程通知中的三
项要求：
\begin{enumerate}
    \item 测量并分析程序运行时间；
    \item 评估算法在图规模增大时的可扩展性；
    \item 将复现方法与至少一个基线方法或变体进行比较。
\end{enumerate}

复现方法是 Kernighan--Lin 细化。两个基线分别是贪心均衡划分和递归谱二分。
贪心方法是快速构造型基线，谱二分则是基于图拉普拉斯特征向量的更强算法基线。

\section{实验设置}

所有实验均通过统一实验入口运行：
\[
\texttt{python src/experiments.py --output-dir outputs --seed 7}.
\]
输出结果保存在 \texttt{outputs/partition\_results.csv} 中。所有方法使用相同的
图实例、分区数、评价指标和随机种子，因此比较是直接且可复现的。

benchmark 包含：
\begin{itemize}
    \item 边长为 6、8、10 的网格图；
    \item 40 和 80 个顶点的 Erdos--Renyi 随机图；
    \item 50 和 90 个顶点的随机几何图；
    \item Zachary karate club 图。
\end{itemize}

大部分合成图使用 $k=4$ 个分区。karate 图较小，并且天然与二分社区结构相关，
因此使用 $k=2$。

\section{评价指标}

主要划分质量指标是割边权重：
\[
    \mathrm{cut}(p)
    =
    \sum_{(u,v)\in E,\;p(u)\neq p(v)} w_{uv}.
\]
割边权重越低，表示跨工作节点 依赖越少。

负载均衡指标为
\[
    \rho(p)=
    \frac{\max_i |\{v:p(v)=i\}|}{|V|/k}.
\]
当值为 1 时表示完全均衡；接近 1 表示算法没有把过多顶点分配给同一个 工作节点。

运行时间以秒为单位，测量计算划分所需的 wall-clock 时间。此外，一致性仿真报告
跨分区消息数量。由于仿真中每条跨工作节点 边在每次迭代会产生两个方向的消息，
因此更低的割值应当对应更少的跨工作节点 通信量。

\section{割边权重对比}

在所有 benchmark 实例上，平均割边权重为：
\[
    \text{greedy}=98.25,\qquad
    \text{spectral}=72.75,\qquad
    \text{Kernighan--Lin}=66.88.
\]
因此，Kernighan--Lin 细化具有最好的平均划分质量。它显著优于贪心基线，并且
平均上也略优于谱划分基线。

\begin{center}
\begin{tabular}{llrrr}
\toprule
图类型 & 规模 & 贪心 & 谱划分 & Kernighan--Lin \\
\midrule
grid & 6 & 23 & 19 & 18 \\
grid & 8 & 40 & 20 & 20 \\
grid & 10 & 65 & 32 & 29 \\
erdos\_renyi & 40 & 54 & 44 & 37 \\
erdos\_renyi & 80 & 122 & 151 & 119 \\
geometric & 50 & 179 & 136 & 136 \\
geometric & 90 & 278 & 155 & 153 \\
karate & 34 & 25 & 25 & 23 \\
\bottomrule
\end{tabular}
\end{center}

结果显示出几个规律。在网格图上，谱二分已经很强，因为图具有清晰的几何结构，
但 Kernighan--Lin 仍然能够进一步改进或持平。在随机几何图上，谱划分和
Kernighan--Lin 都明显优于贪心划分。在较大的 Erdos--Renyi 图上，谱划分反而
比贪心基线更差，说明随机图中的谱结构不一定对应低平衡割；Kernighan--Lin 细化
能够修正初始解并取得更低割值。

\section{运行时间分析}

贪心基线始终是最快方法。这符合预期，因为它只进行一次构造过程，不需要求解
特征值问题，也不需要执行局部搜索细化。

谱二分比贪心更慢，因为它需要计算图拉普拉斯矩阵的特征向量。当前实现使用稠密
特征分解，因此随着图规模增大，成本会快速上升。对于课程规模图，这是可以接受
的；若扩展到更大图，则应使用稀疏特征求解器。

Kernighan--Lin 细化通常最慢，因为它先计算谱划分，再执行额外的成对局部细化。
额外运行时间是为更低割值付出的代价。当通信成本比一次性划分成本更重要时，这
种权衡是合理的，尤其是在划分结果会被许多轮分布式优化迭代重复使用的情况下。

\section{可扩展性评估}

可扩展性可以通过同一图族中不同规模实例进行观察。对于网格图，顶点数从 36 增
加到 64 再到 100，边数从 60 增加到 112 再到 180。所有方法的割值都会随规模
增加而上升，但相对排序比较稳定：贪心割值最高，谱划分明显更低，
Kernighan--Lin 最好或并列最好。

对于 Erdos--Renyi 图，从 40 个顶点增加到 80 个顶点后，图更稠密且更难划分。
较大实例说明单一基线并不总是可靠：谱二分割值为 151，差于贪心的 122，而
Kernighan--Lin 得到 119。这表明当谱方向不是理想切分时，局部细化可以发挥修正
作用。

对于随机几何图，从 50 个顶点增加到 90 个顶点时，边数也显著增加。谱划分依然
表现较强，因为空间结构有意义；Kernighan--Lin 在较大实例上将割值从 155 小幅
改进到 153。这说明当全局几何结构已经被谱方法较好捕捉时，局部细化通常带来
稳定但幅度较小的改进。

\section{通信评估}

一致性仿真报告跨分区消息数。平均跨工作节点 消息数为：
\[
    \text{greedy}=15460,\qquad
    \text{spectral}=11360,\qquad
    \text{Kernighan--Lin}=10440.
\]
这些数字与割边权重趋势一致。这符合预期，因为跨分区消息由端点位于不同分区的
边产生。更低的割边权重对应更低的通信放置成本。

同一张图上，不同划分的最终一致性误差几乎相同。这是因为仿真保持原始图拓扑
不变，仅使用划分来统计消息是否跨越 工作节点边界。这一设计将通信成本与收敛
动力学分开考察。更进一步的仿真可以让跨工作节点 通信变慢、产生延迟或受到预算
限制，从而使划分也影响收敛速度。

\section{结果图}

生成的图像总结了主要实验结果：
\begin{figure}[ht]
    \centering
    \includegraphics[width=0.78\linewidth]{outputs/figures/cut_weight.png}
    \caption{割边权重对比。数值越低，表示跨工作节点 通信越少。}
\end{figure}

\begin{figure}[ht]
    \centering
    \includegraphics[width=0.78\linewidth]{outputs/figures/runtime_seconds.png}
    \caption{运行时间对比。贪心最快；Kernighan--Lin 需要额外细化时间。}
\end{figure}

\begin{figure}[ht]
    \centering
    \includegraphics[width=0.78\linewidth]{outputs/figures/balance_ratio.png}
    \caption{负载均衡比例对比。benchmark 中所有方法都保持了较好均衡性。}
\end{figure}

\section{结论}

实验结果满足本阶段要求：报告测量了运行时间，通过不同规模图观察了可扩展性，
并将 Kernighan--Lin 细化与两个基线进行了比较。结果表明，复现的细化方法通常
能够取得最低割值和最低跨工作节点 通信量，但需要额外运行时间。这支持了项目的
核心观点：经典图划分算法能够帮助更通信高效地组织分布式计算。

\end{document}

